% Source: Wald GR, physical PDF p. 68
%         (printed p. 55).
% Coverage: Chapter 4 opening material before Section 4.1, complete.
% Figures/tables/footnotes: none.
% Status: source-reviewed and compile-clean.
% Uncertainties: none.

In this chapter we shall give a mathematically precise formulation of the ideas
sketched in the introduction. We begin by giving an exposition of the elementary
topic of the geometry of space and the spatial tensorial character of physical laws
in prerelativity physics. The discussion of special relativity which follows will
completely parallel this discussion, with \textquotedblleft
spacetime\textquotedblright{} replacing \textquotedblleft
space\textquotedblright. Our next (and most important) task will be to formulate
general relativity and provide a motivation for Einstein's equation, which relates
the geometry of spacetime to the distribution of matter in the universe. Finally,
we will examine general relativity in the limit where gravity is weak. We will show
that in appropriate circumstances general relativity reproduces the predictions of
the Newtonian theory of gravity. We also will show that in general relativity,
gravity has dynamical degrees of freedom which can be excited by the motion of
matter, i.e., that general relativity predicts the existence of gravitational
radiation.

